\khmerlesson{មេរៀនទី៣៖ ដែននិងកម្លាំងម៉ាញេទិក}
\tanbgdivider

\khmertopic{១. មេដែកនឹងដែនម៉ាញេទិក}
\khmersubtopic{១.១. មេដែក}
\noindent\hspace*{0.6cm}\faCaretRight~\textbf{\textcolor{brownheader}{មេដែក}} គឺជាអង្គធាតុដែលអាចឆក់ទាញដែកបាន។
លក្ខណៈឆក់ទាញទាំងនោះហៅថា \textbf{លក្ខណៈម៉ាញេទិក}
។ ជាទូទៅ មេដែកមានប៉ូលចំនួនពីរគឺ
ចុងម្ខាងជាប៉ូលជើង ({\englishfont N})
និងចុងម្ខាងជាប៉ូលជើង ({\englishfont S})។
ប៉ូលទាំងមិនអាចកាត់ផ្តាច់ពីគ្នាបានឡើយ
ហើយគេចែកវាជា២ប្រភេទគឺ៖

\noindent
\begin{minipage}[t]{0.55\textwidth}
	\raggedright
	\begin{itemize}[label={\textcolor{brownheader}{\faCaretRight}}]
		\item \textbf{មេដែកធម្មជាតិ} ជាមេដែកដែលមានស្រាប់នៅក្នុងធម្មជាតិដូចជា សំណាក \ce{Fe3O4} និងមេដែកផែនដី សញ្ញាជាដើម។
	\begin{itemize}[label={\textcolor{brownheader}{\faCaretRight}}]
		\item \textbf{មេដែកសិប្បនិម្មិត} ជាមេដែកដែលត្រូវបានបង្កើតឡើងតាមតម្រូវការចាំបាច់ មានដូចជា របាមេដែក ម្ជុលមេដែក និងមេដែករាងអក្សរ ({\englishfont U})។
	\end{itemize}
	\end{itemize}
	\khmersubsubtopic{ក. អន្តរកម្មនៃមេដែក}
	នៅពេលដែលគេដាក់មេដែកពីរជិតគ្នា
	នោះបាតុភូតពីរនឹងកើតឡើង
	នោះគឺមេដែកច្រានចេញពីគ្នា ឬទាញគ្នាចូល
	ដែលបាតុភូតនេះហៅថា
	\textbf{អន្តរកម្មនៃមេដែក ឬអន្តរកម្មមេដែក}
	និងមានលក្ខណៈដូចខាងក្រោម៖
	\begin{itemize}[label={\textcolor{brownheader}{\faCaretRight}}]
		\item បើគេដាក់ប៉ូលដូចគ្នា ({\englishfont N-N} ឬ {\englishfont S-S}) ជិតគ្នា នោះមេដែកទាំងពីរនឹងច្រានចេញពីគ្នា។
		\item បើគេដាក់ប៉ូលខុសគ្នា ({\englishfont N-S} ឬ {\englishfont S-N}) ជិតគ្នា នោះមេដែកទាំងពីរនឹងទាញគ្នាចូល។
	\end{itemize}
	\end{minipage}%
\hfill
\begin{minipage}[t]{0.40\textwidth}
	\vspace{0pt}
	\khmerphoto[0.95\linewidth]{img/magnetite.jpg}{សំណាកមេដែក \ce{Fe3O4}}
	\vspace{0.75mm}
	\khmerphoto[0.75\linewidth]{img/mgshapes.jpg}{រាងមេដែកសិប្បនិម្មិត}
\end{minipage}

\begin{center}
	\newcommand{\drawmagnet}[3]{%
		\begin{scope}[shift={(#1,#2)}]
		\ifnum#3=0
		\fill[waveblue!80] (-1.05,-0.30) rectangle (0,0.30);
		\fill[wavered!80] (0,-0.30) rectangle (1.05,0.30);
		\node[white, font=\bfseries\small] at (-0.525,0) {S};
		\node[white, font=\bfseries\small] at (0.525,0) {N};
		\else
			\fill[wavered!80] (-1.05,-0.30) rectangle (0,0.30);
			\fill[waveblue!80] (0,-0.30) rectangle (1.05,0.30);
			\node[white, font=\bfseries\small] at (-0.525,0) {N};
			\node[white, font=\bfseries\small] at (0.525,0) {S};
		\fi
		\fill[white, opacity=0.15] (-1.05,0.22) rectangle (1.05,0.30);
		\fill[black, opacity=0.10] (-1.05,-0.30) rectangle (1.05,-0.22);
		\draw[black!20, line width=0.35pt, rounded corners=0.6pt] (-1.05,-0.30) rectangle (1.05,0.30);
		\end{scope}%
	}
	
	\begin{center}
		\begin{tikzpicture}[
			font=\small,
			>=Stealth,
			line join=round,
			line cap=round,
			scale=1.05
		]

% --- Row 1: Attract (N faces S) ---
			\drawmagnet{-2.0}{1.6}{0}   % left: S|N
			\drawmagnet{2.0}{1.6}{0}    % right: S|N
			
			\draw[black, line width=1.5pt, ->] (-0.85,1.6) -- (-0.15,1.6);
			\draw[black, line width=1.5pt, ->] (0.85,1.6) -- (0.15,1.6);
			\node[above=7pt, brownheader] at (
			0,1.6) {ទាញចូល};

% --- Row 2: Repel (N faces N) ---
			\drawmagnet{-2.0}{0}{0}     % left: S|N
			\drawmagnet{2.0}{0}{1}      % right: N|S
			
			\draw[black, line width=1.5pt, ->] (-0.15,0) -- (-0.85,0);
			\draw[black, line width=1.5pt, ->] (0.15,0) -- (0.85,0);
			\node[above=7pt, brownheader] at (
			0,0) {ច្រានចេញ};

% --- Row 3: Repel (S faces S) ---
			\drawmagnet{-2.0}{-1.6}{1}  % left: N|S
			\drawmagnet{2.0}{-1.6}{0}   % right: S|N
			
			\draw[black, line width=1.5pt, ->] (-0.15,-1.6) -- (-0.85,-1.6);
			\draw[black, line width=1.5pt, ->] (0.15,-1.6) -- (0.85,-1.6);
			\node[above=7pt, brownheader] at (
			0,-1.6) {ច្រានចេញ};
		
		\end{tikzpicture}
	\end{center}
	\textcolor{brownheader}{\khmerpagenum អន្តរកម្មនៃមេដែក}
\end{center}

\khmersubtopic{១. ២. ដែនម៉ាញេទិកនៃមេដែក}
លក្ខណៈទាញគ្នាចូលនិងច្រានចេញនៃមេដែកគឺបានបង្កើតឡើងដោយសារកម្លាំងអន្តរកម្មម៉ាញេទិកនៃមេដែកទាំងនោះ។ កាលពីថ្នាក់ទី១១ យើងបានសិក្សារួចមកហើយថា កម្លាំងអគ្គិសនីដែលមានអំពើលើបន្ទុកអគ្គិសនីត្រូវបានពណ៌នាដោយលំហនៃដែនអគ្គិសនីជុំវិញបន្ទុកអគ្គិសនីស្ថិតនៅនឹង ហើយបង្កើតឲ្យមានអន្តរកម្មរវាងដែនអគ្គិសនីនិងបន្ទុកអគ្គិសនីនោះ។ នៅក្នុងការសិក្សាអំពីកម្លាំងម៉ាញេទិកក៏មានលក្ខណៈស្រដៀងទៅនឹងកម្លាំងអគ្គិសនីដែរ គឺមានន័យថាកម្លាំងម៉ាញេទិកដែលត្រូវបានបង្កើតឡើងដោយលំហនៃដែនម៉ាញេទិក មានអន្តរកម្មលើមេដែកដែលស្ថិតនៅក្នុងលំហនោះ។ ក្នុងការសិក្សាដែនម៉ាញេទិក គ៏មានលក្ខណៈស្រដៀងទៅនឹងកម្លាំងអគ្គិសនីដែរ គឺមានន័យថាកម្លាំងម៉ាញេទិក ដែលបង្កើតឡើងនៅក្នុងលំហនៃដែនម៉ាញេទិកមានអន្តរកម្មលើមេដែកស្ថិតនៅក្នងលំហនោះ។ ដូចនេះ ដែនម៉ាញេទិក ($\vec{B}$) គឺជាមជ្ឍដ្ឋានរូបធាតុពិសេសសដែលកើតឡើងជុំវិញមេដែក ហើយនឹងបង្កើតកម្លាំងម៉ាញេទិកលើមេដែក (ច្រានចេញ និងទាញចូល) ស្ថិតនៅក្នុងមជ្ឈដ្ឋានមួយនេះ។\\

\noindent
\begin{minipage}[t]{0.55\textwidth}
	\raggedright
ដើម្បីបង្ហាញពីលក្ខណៈវុិចទ័រនៃដែនម៉ាញេទិក
គេត្រូវធ្វើការពិសោធន៍ដោយរោយកម្ទេចដែក
នៅលើក្រដាសរឹងមួយសន្លឹក នៅជុំវិញរបារមេដែកមួយ
បន្ទាប់មកគេគោះក្រដាសនោះតិចៗ
ហើយយើងសង្កេតឃើញថាកម្ទេចដែកនោះត្រូវបានរត់ជាជួរខ្សែមានរាយជុំវិញមេដែក។ កម្ទេចដែលបានរៀបជាជួរហៅថា ស្បុិចម៉ាញេទិកនៃរបារមេដែក។\\
\end{minipage}%
\hfill
\begin{minipage}[t]{0.45\textwidth}
	\vspace{0pt}
	\khmerphoto[0.5\linewidth]{
img/magfield.jpg}{
ស្បុិចម៉ាញេទិកនៃរបារមេដែក}
\end{minipage}

\noindent
\begin{minipage}[t]{0.55\textwidth}
	\raggedright
	\faCaretRight~\textbf{\textcolor{brownheader}{ខ្សែដែនម៉ាញេទិក}} គឺជាខ្សែដែលប៉ះនឹងវុិចទ័រដែនម៉ាញេទិក $\vec{B}$ ហើយនៅគ្រប់ចំណុចនីមួយៗរបស់វាហើយមានវុិចទ័រដែនម៉ាញេទិក $\vec{B}$។\\
	លក្ខណៈនៃខ្សែដែនម៉ាញេទិកត្រូវបានកំណត់ដោយ៖
	\begin{itemize}[label={\textcolor{brownheader}{\faCaretRight}}]
	\item ចាប់ផ្តើមពីប៉ូលជើង {\englishfont N}
	ហើយបញ្ចប់នៅប៉ូលត្បូង {\englishfont S}។
	\item មានទឹសដៅពីលើប៉ូលជើង {\englishfont N}
	ទៅប៉ូលត្បូង {\englishfont S} ផ្នែកខាងក្នុង។
	\item មានរាងជាខ្សែបិទនិងមិនកាត់គ្នាទេ។
	\end{itemize}
\end{minipage}%
\hfill
\begin{minipage}[t]{0.40\textwidth}
	\vspace{0pt}
	\centering
	\begin{tikzpicture}[
		font=\small,
		>={Stealth[scale=0.55, length=2.2mm, width=1.8mm]},
		line join=round,
		line cap=round,
		scale=0.85
	]

% --- Parameters ---
		\def\mw{1.9}
		\def\mh{0.38}
		\def\ymax{2.4}

% ==================== SUBTLE BACKGROUND GLOW ====================
		\fill[primary, opacity=0.04] (0,0) ellipse (3.6cm and 2.7cm);

% ==================== MAGNET BODY ====================
		\fill[wavered!72] (-\mw,-\mh) rectangle (0,\mh);
		\fill[waveblue!72] (0,-\mh) rectangle (\mw,\mh);

% Top highlight / bottom shadow
		\fill[white, opacity=0.18] (-\mw,\mh-0.04) rectangle (\mw,\mh);
		\fill[black, opacity=0.12] (-\mw,-\mh) rectangle (\mw,-\mh+0.04);

% Soft border
		\draw[black!14, line width=0.35pt, rounded corners=0.6pt] (-\mw,-\mh) rectangle (\mw,\mh);
		\draw[black!10, line width=0.2pt] (0,-\mh) -- (0,\mh);

% Labels
		\node[white, font=\bfseries\small] at (-\mw/2,0) {\englishfont N};
		\node[white, font=\bfseries\small] at (\mw/2,0) {\englishfont S};

% ==================== INTERNAL FIELD LINES ====================
		\foreach \y in {-0.12, 0, 0.12} {
			\draw[primary!28, dashed, line width=0.45pt, ->] (\mw-0.10,\y) -- (-\mw+0.10,\y);
		}

% ==================== EXTERNAL FIELD LINES ====================
% Loop 1 (innermost)
		\draw[primary!78, line width=0.85pt] (-\mw+0.02,\mh+0.01) to[out=68, in=180] (0,0.52);
		\draw[primary!78, line width=0.85pt, <-] (0,0.52) to[out=0, in=112] (\mw-0.02,\mh+0.01);
		\draw[primary!78, line width=0.85pt] (-\mw+0.02,-\mh-0.01) to[out=-68, in=180] (0,-0.52);
		\draw[primary!78, line width=0.85pt, <-] (0,-0.52) to[out=0, in=-112] (\mw-0.02,-\mh-0.01);

% Loop 2
		\draw[primary!65, line width=0.75pt] (-\mw+0.02,\mh+0.01) to[out=78, in=180] (0,0.98);
		\draw[primary!65, line width=0.75pt, <-] (0,0.98) to[out=0, in=102] (\mw-0.02,\mh+0.01);
		\draw[primary!65, line width=0.75pt] (-\mw+0.02,-\mh-0.01) to[out=-78, in=180] (0,-0.98);
		\draw[primary!65, line width=0.75pt, <-] (0,-0.98) to[out=0, in=-102] (\mw-0.02,-\mh-0.01);

% Loop 3
		\draw[primary!52, line width=0.68pt] (-\mw+0.02,\mh+0.01) to[out=84, in=180] (0,1.48);
		\draw[primary!52, line width=0.68pt, <-] (0,1.48) to[out=0, in=96] (\mw-0.02,\mh+0.01);
		\draw[primary!52, line width=0.68pt] (-\mw+0.02,-\mh-0.01) to[out=-84, in=180] (0,-1.48);
		\draw[primary!52, line width=0.68pt, <-] (0,-1.48) to[out=0, in=-96] (\mw-0.02,-\mh-0.01);

% Loop 4 (outermost)
		\draw[primary!40, line width=0.6pt] (-\mw+0.02,\mh+0.01) to[out=88, in=180] (0,2.05);
		\draw[primary!40, line width=0.6pt, <-] (0,2.05) to[out=0, in=92] (\mw-0.02,\mh+0.01);
		\draw[primary!40, line width=0.6pt] (-\mw+0.02,-\mh-0.01) to[out=-88, in=180] (0,-2.05);
		\draw[primary!40, line width=0.6pt, <-] (0,-2.05) to[out=0, in=-92] (\mw-0.02,-\mh-0.01);

% ==================== FAR-FIELD SIDE LINES ====================
		\draw[primary!35, line width=0.55pt, ->] (-\mw-0.03, 0.14) -- (-\mw-1.6, 0.14);
		\draw[primary!35, line width=0.55pt, ->] (-\mw-0.03,-0.14) -- (-\mw-1.6,-0.14);
		\draw[primary!35, line width=0.55pt, ->] (\mw+1.6, 0.14) -- (\mw+0.03, 0.14);
		\draw[primary!35, line width=0.55pt, ->] (\mw+1.6,-0.14) -- (\mw+0.03,-0.14);
	
	\end{tikzpicture}
	
	\begin{center}
		\textcolor{brownheader}{\khmerpagenum ខ្សែដែនម៉ាញេទិកនៃរបារមេដែក}
	\end{center}
	
\end{minipage}

\khmersubtopic{១. ៣. ដែនម៉ាញេទិកនៃផែនដី}
យើងបានយល់ច្បាស់អំពីអន្តរកម្មរវាងមេដែកពីរ
នៅពេលដែលដាក់មេដែកដែលមានប៉ូលដូចនោះវានឹងច្រានគ្នាចេញ
ឬមេដែកដែលមានប៉ូលខុសគ្នានោះវានឹងទាញគ្នាចូល។
នៅពេលនេះ យើងនឹងធ្វើការសួរថា "\textbf{ហេតុអ្វីបានជាប៉ូល {\englishfont N} នៃម្ជុលមេដែករបស់ត្រីវិស័យចង្អុលទៅកាន់ទឹសខាងជើងនៃប៉ូលរបស់ផែនដី? ឬហេតុអ្វីបានជារបារមេដែកដែលត្រូវបានគេចងភ្ជាប់ និងព្យួរវាដោយខ្សែអំបោះត្រង់ពាក់កណ្តាល (ទីប្រជុំទម្ងន់) បន្ទាប់មកទុកឲ្យនៅនឹងដោយខ្លួនឯងស្រាប់តែប៉ូលជើងនៃរបារមេដែកនោះចង្អុលទៅកាន់ប៉ូលជើងនៃផែនដីទៅវិញ?}"\\
\noindent
\begin{minipage}[t]{0.45\textwidth}
	\raggedright
	អ្នកវិទ្យាសាស្ត្របានបង្ហាញថា
	ដោយសារតែនៅក្នុងស្នូលផែនដីសម្បូរទៅដោយលោហៈនិងផែនដីបានធ្វើចលនារង្វិលជុំវីញអ័ក្សរបស់វា
	ដែលបង្កើតបានជាចរន្តអគ្គិសនីផ្លាស់ទីនៅលើគន្លងវង់ និងមានដែនម៉ាញេទិកស្ថិតនៅចំកណ្តាលផែនដី ដែលមានអ័ក្សមួយដែលពុំកាត់អ័ក្សរង្វិលរបស់ផែនដី ប្រៀបបានដូចជារបារមេដែកកប់ជ្រៅនៅក្នុងផែនដី។
	របារមេដែកនៃផែនដីមានប៉ូលជើង {\englishfont N}
	ម៉ាញេទិកនៃរបារមេដែកផែនដីស្ថិតនៅជិតនឹងប៉ូលត្បូងនៃផែនដីភូមិសាស្ត្រនៃផែនដី ហើយមានប៉ូលត្បូង {\englishfont S}
	ភូមិសាស្ត្រនៃផែនដី។
\end{minipage}%
\hfill
\begin{minipage}[t]{0.55\textwidth}
	\vspace{0pt}
	\khmerphoto{
		img/earthmgfield.jpg}{
		ស្បុិចម៉ាញេទិកនៃរបារមេដែកផែនដី}
\end{minipage}
\noindent
\begin{minipage}[t]{0.55\textwidth}
	\raggedright
	\begin{itemize}[label={\textcolor{brownheader}{\faCaretRight}}]
    \item ម្ជុលមេដែកមិនចង្អុលទៅប៉ូលជើងពិតនោះទេ(
    លើកលែងតែនៅតំបន់ជិតខ្សែអេក្វាទ័រ)
    ព្រោះប៉ូលជើងនៃម៉ាញេទិកនៃផែនដី
    និងប៉ូលជើងនៃភូមិសាស្ត្រ ស្ថិតនៅក្នុងទីតាំងពីរផ្សេងគ្នា(មិនត្រួតស៊ីគ្នារវាងអ័ក្សរង្វិលនៃផែនដី និងអ័ក្សនៃមេដែកនៃផែនដី) បើគិតរង្វាស់មុំដែលមានគម្លាតពីគ្នានៅត្រង់ចំណុចប្រសព្វរវាងអ័ក្សរង្វិលនៃផែនដី និងអ័ក្សនៃចរន្តមេដែកនៃផែនដីគឺមានតម្លៃប្រហែល $15^\circ$។ តម្លៃនេះត្រូវបានគេហៅថា \textbf{ដេគ្លីនេស្យុង (declination)} ។
    \item ប្រសិនជាគេព្យួរម្ជុលមេដែកនៃទ្រនិចនាឡិកាទៅនឹងអ័ក្សរង្វិលស្ថិតនៅក្នុងប្លង់ដេក និងឲ្យទ្រនិចនាឡិកាអាចរំកិលបានក្នុងប្លង់ឈរបានដោយសេរី នោះនៅពេលដែលមានលំនឹង គេសង្កេតឃើញថាប៉ូល S ងាកទៅក្រោមធៀបនឹងប្លង់ដេក រឺប៉ូល N ងើបឡើងលើបានមុំមួយ។ មុំដែលធៀបនឹងប្លង់ដេកនេះត្រូវបានគេហៅថា \textbf{អ៊ីគ្លីនេស្យុង (inclination)} ។
	\end{itemize}
\end{minipage}%
\hfill
\begin{minipage}[t]{0.45\textwidth}
	\vspace{0pt}
	\khmerphoto[0.5\linewidth]{
		img/declination.jpg}{
		រង្វាស់មុំដេគ្លីនេស្យុងនៃផែនដី}
	\khmerphoto[0.5\linewidth]{
		img/inclination.png}{
		រង្វាស់មុំអ៊ីគ្លីនេស្យុងនៃផែនដី}
\end{minipage}

\khmertopic{២. ដែនម៉ាញេទិកដែលបង្កើតឡើងដោយចរន្តអគ្គិសនី}
\khmersubtopic{២.១. ដែនម៉ាញេទិកនៃចរន្តត្រង់}

\noindent
\begin{minipage}[t]{0.55\textwidth}
\raggedright
\faChevronRight\ \textcolor{brownheader}{\textbf{ពិសោធន៍៖ }}
យកខ្សែចម្លងត្រង់មួយ
ដោតចូលប្រហោងនៃផង់កាតុងដែលបានចោះរន្ធចំកណ្តាលនៃផ្ទៃក្រដាស បន្ទាប់មកយកកម្ទេចដែកមក រោយលើក្រដាសជុំវិញខ្សែចម្លង។ បន្ទាប់យើងផ្តល់ចរន្ត $I$ ឲ្យរត់នៅក្នុងខ្សែចម្លង ហើយគោះក្រដាសតិចៗ នៅពេលនោះយើងសង្កេតឃើញថាកម្ទេចដែកតម្រៀបគ្នាតាមរាងជារង្វង់ជុំវិញខ្សែចម្លង។ កម្ទេចដែកដែលតម្រៀបគ្នានោះគឺជាខ្សែចរនៃដែនម៉ាញេទិក តាងឲ្យខ្សែដែនម៉ាញេទិចដែលបង្កើតឡើងដោយខ្សែចម្លងត្រង់ ។
\end{minipage}%
\hfill
\begin{minipage}[t]{0.40\textwidth}
	\vspace{0pt}
	\hspace*{0.5em}%
	\khmerphoto[0.95\linewidth]{
		img/straightwirefield.png}{
		ខ្សែដែនម៉ាញេទិចដែលបង្កើតឡើងដោយខ្សែចម្លងត្រង់}
\end{minipage}

\vspace{1em}
\noindent
\begin{minipage}[t]{0.55\textwidth}
	\raggedright
	\textbf{លក្ខណៈនៃវ៉ិចទ័រដែនម៉ាញ៉េទិក (
		$\vec{B}$) បង្កើតដោយចរន្តត្រង់នៅត្រង់ចំណុច M ៖}

	\vspace{0.5em}
	\begin{itemize}[itemsep=4pt, topsep=2pt, parsep=0pt, leftmargin=1.6em]
	\item \textbf{ចំណុចចាប់} ៖ ត្រង់ចំណុច M ។
	\item \textbf{ទិស} ៖ កែងនឹងប្លង់ខ្សែចម្លងដែលកំណត់ដោយ $PP'$ និង $OM$ ។
	\item \textbf{ទិសដៅ} ៖ ចេញក្រៅពីផ្ទៃបាតដៃ (
	វិធានដៃស្តាំ ៖
	ដាក់មេដៃស្របនឹងទិសដៅចរន្តអគ្គិសនីរត់ក្នុងខ្សែចម្លង ម្រាមទាំង៤កោងចង្អុលចំណុចបង្កើតដែនម៉ាញ៉េទិក M បន្ទាប់មកក្ដាប់បង្ហាញពីទិសដៅដែនម៉ាញ៉េទិកក្រង់ចំណុចនោះ) ។
	\item \textbf{តម្លៃ} ៖ $B = \mu_0 \dfrac{I}{2\pi d}$ \quad (១)
	\end{itemize}
	\vspace{0.5em}
	\defbox{
	\begin{center}
		$B$ ៖ អាំងតង់ស៊ីតេដែនម៉ាញ៉េទិក ($\mathrm{T}$)\\[1pt]
		$\mu_0 = 4\pi \times 10^{-7}\ \mathrm{T\cdot m/A}$ ៖ ជម្រាបម៉ាញ៉េទិកនៅក្នុងខ្យល់ ឬ សូន្យកាស\\[1pt]
		$I$ ៖ ចរន្តអគ្គិសនី ($\mathrm{A}$)
	\end{center}
	}
\end{minipage}%
\hfill
\begin{minipage}[t]{0.40\textwidth}
	\vspace{0pt}
	\hspace*{0.5em}%
	\khmerphoto[0.95\linewidth]{
		img/righthandgriprule.jpg}{
		វិធានដៃស្តាំដើម្បីរកទិសដៅដែនម៉ាញ៉េទិកបង្កើតឡើងដោយខ្សែចម្លងត្រង់}
\end{minipage}

\begin{center}
	\textbf{ចំណាំ៖  }
	លក្ខខណ្ឌរូបមន្តខាងលើអាចប្រើបានលុះត្រាតែប្រវែងខ្សែចម្លង ($l$)
	មានប្រវែងអានន្តតែប៉ុណ្ណោះ។
\end{center}

\khmersubtopic{២.២. ដែនម៉ាញេទិកនៃចរន្តវង់}

\noindent
\begin{minipage}[t]{0.55\textwidth}
	\raggedright
	\faChevronRight\ \textcolor{brownheader}{\textbf{ពិសោធន៍៖ }}
	យើងយកខ្សែចម្លងរាងជារង្វង់មកដោតចូលក្នុងរន្ធទាំងពីរនៃប្លង់ក្រដាស ដោយលៃយ៉ាងណាកន្លះរង្វង់នៅផ្នែកខាងលើ
	និងកន្លះរង្វង់នៅផ្នែកខាងក្រោម
	បន្ទាប់មករោយកម្ទេចដែកនៅលើផ្ទៃក្រដាស
	ហើយឲ្យចរន្តរត់ឆ្លងកាត់ខ្សែចម្លងនោះ ។
	យើងគោះក្រដាសនោះតិចៗ
	បន្ទាប់មកកម្ទេចដែកតម្រៀបគ្នាបង្កើតបានជា
	ស្បុិចម៉ាញ៉េទិក
	ដែលត្រូវបានគេយកគម្រូនៃស៊ិចនោះមកគូសបានជា
	ខ្សែដែនម៉ាញ៉េទិក។

\end{minipage}%
\hfill
\begin{minipage}[t]{0.40\textwidth}
	\vspace{0pt}
	\hspace*{0.5em}%
	\khmerphoto[0.95\linewidth]{img/roundedwirefield.jpg}{ខ្សែដែនម៉ាញេទិចដែលបង្កើតឡើងដោយខ្សែចម្លងវង់}
\end{minipage}

\vspace{1em}
\noindent
\begin{minipage}[t]{0.55\textwidth}
	\raggedright
	\textit{\textgreater\
	\textbf{លក្ខណៈនៃវ៉ិចទ័រដែនម៉ាញ៉េទិក ($\vec{B}$) បង្កើតដោយចរន្តវង់នៅត្រង់ផ្ចិតនៃរង្វង់ ៖}}

	\vspace{0.5em}
	\begin{itemize}[itemsep=4pt, topsep=2pt, parsep=0pt, leftmargin=1.6em]
	\item \textbf{ចំណុចចាប់} ៖ ត្រង់ផ្ចិតនៃរង្វង់
	\item \textbf{ទិស} ៖ កែងនឹងប្លង់នៃខ្សែចម្លង
	\item \textbf{ទិសដៅ} ៖ វិធានបាតដៃស្តាំ (
	ដាក់បាតដៃក្តោបលើស្នៀតវង់ លៃយ៉ាងណាឲ្យចរន្ត ចេញតាមម្រាមទាំងបួន រីឯមេដៃកែងបង្ហាញពីទិសដៅដែនម៉ាញ៉េទិក) ។
	\item \textbf{តម្លៃ} ៖ $B = \mu_0 \dfrac{I}{2R}$ \quad (២)
	\end{itemize}
	\vspace{0.5em}
	\defbox{
	\begin{center}
		$B$ ៖ អាំងតង់ស៊ីតេដែនម៉ាញ៉េទិកបង្កើតដោយចរន្តវង់ ($\mathrm{T}$)\\[1pt]
		$\mu_0 = 4\pi \times 10^{-7}\ \mathrm{T\cdot m/A}$ ៖ ជម្រាបម៉ាញ៉េទិក\\[1pt]
		$I$ ៖ ចរន្តអគ្គិសនី ($\mathrm{A}$)\\[1pt]
		$R$ ៖ កាំនៃបុំខ្សែចម្លងរាងវង់ ($\mathrm{m}$)
	\end{center}
	}
\end{minipage}%
\hfill
\begin{minipage}[t]{0.40\textwidth}
	\vspace{0pt}
	\hspace*{0.5em}%
	\khmerphoto[0.95\linewidth]{
		img/righthandroundedrule.jpg}{
		វិធានបាតដៃស្តាំដើម្បីរកទិសដៅដែនម៉ាញ៉េទិចបង្កើតដោយចរន្តវង់}
\end{minipage}
\begin{center}
	\textbf{ចំណាំ៖  }
	លក្ខខណ្ឌរូបមន្តខាងលើអាចប្រើបានលុះត្រាតែប្រវែងខ្សែចម្លងមិនរុំជាស្ពៀឡើយ។
	តែប្រសិនបើខ្សែចម្លងរុំជាស្ពៀចំនួន $N$
	ស្ពៀនោះ រូបមន្ត (២) នឹងក្លាយទៅជា
	$\boxed{B = \mu_0 \dfrac{NI}{2R}}$។ រូបមន្តរក $N$ នឹងបង្ហាញនៅចំណុចក្រោយ។
\end{center}

\khmersubtopic{២.៣. ដែនម៉ាញេទិកនៃសូលេណូអុីត}

\noindent
\begin{minipage}[t]{0.55\textwidth}
	\raggedright
	\faChevronRight\ \textcolor{brownheader}{\textbf{ពិសោធន៍៖ }}
	យើងយកខ្សែចម្លងមកដោតចូលក្នុងរន្ធនៃប្លង់ក្រដាស បង្កើតបានជារាងសូលេណូអ៊ីត(ឬប៊ីនសំរ៉ែត) ដ
	ោយលៃយ៉ាងណាកន្លះរង្វង់នៅផ្នែកខាងលើ និងកន្លះរង្វង់នៅផ្នែកខាងក្រោម បន្ទាប់មករោយកម្ទេចដែកនៅលើផ្ទៃក្រដាស ហ
	ើយឲ្យចរន្តរត់ឆ្លងកាត់ខ្សែចម្លងនោះ ។ យើងគោះក្រដាសនោះតិចៗ បន្ទាប់មកកម្ទេចដែកតម្រៀបគ្នាបង្កើតបានជាស្បុិចម៉ាញ៉េទិក ដ
	ែលត្រូវបានគេយកគម្រូនៃស៊ិចនោះមកគូសបានជា \textbf{\textit{ខ្សែដែនម៉ាញ៉េទិក}}។
	
\end{minipage}%
\hfill
\begin{minipage}[t]{0.40\textwidth}
	\vspace{0pt}
	\hspace*{0.5em}%
	\khmerphoto[0.95\linewidth]{img/solenoid.jpg}{ខ្សែដែនម៉ាញេទិចដែលបង្កើតឡើងដោយសូលេណូអុីត}
\end{minipage}

